% !TEX root = ../main.tex

\chapter{Hasil Pengujian}

Bab ini merangkum rancangan eksperimen dan format evaluasi sesuai spesifikasi tugas. Nilai numerik pada tabel diisi dari hasil notebook pelatihan/evaluasi di \texttt{src/notebooks}.

\section{CNN}

\subsection{Setup Eksperimen}

\begin{itemize}
	\item Dataset: Intel Image Classification (6 kelas).
	\item Loss: Sparse Categorical Crossentropy.
	\item Optimizer: Adam.
	\item Metrik utama: macro F1-score pada split test.
	\item Komparasi utama: shared vs non-shared parameter dan Keras vs from scratch.
\end{itemize}

\subsection{Perbandingan Shared vs Non-Shared Parameter}

\begin{table}[H]
	\centering
	\caption{Perbandingan Conv2D (Shared) vs LocallyConnected2D (Non-Shared)}
	\begin{tabular}{lccc}
		\toprule
		\textbf{Model} & \textbf{Macro F1 (Test)} & \textbf{Jumlah Parameter} & \textbf{Waktu Inferensi} \\
		\midrule
		Conv2D (Keras) & [isi] & [isi] & [isi] \\
		Conv2D (Scratch) & [isi] & [isi] & [isi] \\
		LocallyConnected2D (Keras) & [isi] & [isi] & [isi] \\
		LocallyConnected2D (Scratch) & [isi] & [isi] & [isi] \\
		\bottomrule
	\end{tabular}
	\label{tab:cnn-shared-vs-nonshared}
\end{table}

\textbf{Grafik.} Sertakan grafik training loss, validation loss, dan macro F1 untuk perbandingan Conv2D vs LocallyConnected2D.

\textbf{Analisis.} Jelaskan dampak \textit{parameter sharing} terhadap akurasi, efisiensi parameter, dan biaya komputasi.

\subsection{Pengaruh Jumlah Layer Konvolusi}

\begin{table}[H]
	\centering
	\caption{Pengaruh Jumlah Layer Konvolusi}
	\begin{tabular}{lcc}
		\toprule
		\textbf{Variasi} & \textbf{Macro F1 (Test)} & \textbf{Catatan} \\
		\midrule
		1 layer konvolusi & [isi] & [isi] \\
		2 layer konvolusi & [isi] & [isi] \\
		3 layer konvolusi & [isi] & [isi] \\
		\bottomrule
	\end{tabular}
	\label{tab:cnn-var-depth}
\end{table}

\textbf{Grafik.} Sertakan: (1) kurva macro F1 vs jumlah layer, (2) training loss untuk setiap variasi depth.

\textbf{Analisis.} Diskusikan pengaruh kedalaman model terhadap kapabilitas ekstraksi fitur dan overfitting.

\subsection{Pengaruh Banyak Filter per Layer}

\begin{table}[H]
	\centering
	\caption{Pengaruh Kombinasi Jumlah Filter}
	\begin{tabular}{lcc}
		\toprule
		\textbf{Kombinasi Filter} & \textbf{Macro F1 (Test)} & \textbf{Catatan} \\
		\midrule
		Kombinasi A & [isi] & [isi] \\
		Kombinasi B & [isi] & [isi] \\
		Kombinasi C & [isi] & [isi] \\
		\bottomrule
	\end{tabular}
	\label{tab:cnn-var-filter-count}
\end{table}

\textbf{Grafik.} Sertakan: (1) kurva macro F1 untuk setiap kombinasi filter, (2) plot jumlah parameter vs macro F1.

\textbf{Analisis.} Analisis trade-off antara kapasitas model (jumlah filter) terhadap akurasi dan computational cost.

\subsection{Pengaruh Ukuran Filter per Layer}

\begin{table}[H]
	\centering
	\caption{Pengaruh Ukuran Kernel Konvolusi}
	\begin{tabular}{lcc}
		\toprule
		\textbf{Kombinasi Kernel} & \textbf{Macro F1 (Test)} & \textbf{Catatan} \\
		\midrule
		Kombinasi A & [isi] & [isi] \\
		Kombinasi B & [isi] & [isi] \\
		Kombinasi C & [isi] & [isi] \\
		\bottomrule
	\end{tabular}
	\label{tab:cnn-var-kernel-size}
\end{table}

\textbf{Grafik.} Sertakan: (1) kurva macro F1 untuk setiap kombinasi kernel, (2) comparison spatial receptive field vs macro F1.

\textbf{Analisis.} Diskusikan pengaruh receptive field (ukuran kernel) terhadap kemampuan menangkap pola visual dan akurasi klasifikasi.

\subsection{Pengaruh Jenis Pooling Layer}

\begin{table}[H]
	\centering
	\caption{Perbandingan Jenis Pooling}
	\begin{tabular}{lcc}
		\toprule
		\textbf{Pooling} & \textbf{Macro F1 (Test)} & \textbf{Catatan} \\
		\midrule
		MaxPooling & [isi] & [isi] \\
		AvgPooling & [isi] & [isi] \\
		\bottomrule
	\end{tabular}
	\label{tab:cnn-var-pooling}
\end{table}

\textbf{Grafik.} Sertakan: training loss dan validation loss untuk kedua pooling strategy.

\textbf{Analisis.} Bandingkan efektivitas max pooling vs average pooling dalam menjaga informasi penting dan mengurangi overfitting.

\subsection{Perbandingan Keras vs From Scratch (Arsitektur Terbaik)}

\begin{table}[H]
	\centering
	\caption{Keras vs From Scratch pada Arsitektur CNN Terbaik}
	\begin{tabular}{lcc}
		\toprule
		\textbf{Implementasi} & \textbf{Macro F1 (Test)} & \textbf{Selisih terhadap Keras} \\
		\midrule
		Keras & [isi] & -- \\
		From Scratch & [isi] & [isi] \\
		\bottomrule
	\end{tabular}
	\label{tab:cnn-keras-vs-scratch}
\end{table}

\textbf{Grafik.} Sertakan: (1) kurva training loss (Keras vs Scratch), (2) kurva validation accuracy, (3) prediksi test accuracy.

\textbf{Analisis.} Verifikasi keselarasan numerik implementasi \textit{from scratch} dengan Keras, termasuk backward propagation accuracy jika tersedia.

\section{Simple RNN dan LSTM untuk Image Captioning}

\subsection{Setup Eksperimen}

\begin{itemize}
	\item Dataset: Flickr8k (6000 train, 1000 val, 1000 test).
	\item Encoder: CNN pretrained (frozen), fitur disimpan \texttt{.npy}.
	\item Decoder: RNN dan LSTM, dilatih terpisah.
	\item Loss: Sparse Categorical Crossentropy; optimizer: Adam.
	\item Metrik: BLEU-4 (utama), METEOR (sekunder), dan waktu eksekusi.
\end{itemize}

\subsection{Perbandingan Jumlah Layer: RNN}

\begin{table}[H]
	\centering
	\caption{Variasi Jumlah Layer Recurrent pada Decoder RNN}
	\begin{tabular}{lccc}
		\toprule
		\textbf{Jumlah Layer} & \textbf{BLEU-4} & \textbf{METEOR} & \textbf{Waktu Inferensi} \\
		\midrule
		1 layer & [isi] & [isi] & [isi] \\
		2 layer & [isi] & [isi] & [isi] \\
		3 layer & [isi] & [isi] & [isi] \\
		\bottomrule
	\end{tabular}
	\label{tab:rnn-var-depth}
\end{table}

\textbf{Grafik.} Training loss, validation loss, dan BLEU-4 vs jumlah layer.

\textbf{Analisis.} Diskusikan pengaruh kedalaman decoder RNN terhadap representasi kontekstual dan kualitas caption.

\subsection{Perbandingan Jumlah Layer: LSTM}

\begin{table}[H]
	\centering
	\caption{Variasi Jumlah Layer Recurrent pada Decoder LSTM}
	\begin{tabular}{lccc}
		\toprule
		\textbf{Jumlah Layer} & \textbf{BLEU-4} & \textbf{METEOR} & \textbf{Waktu Inferensi} \\
		\midrule
		1 layer & [isi] & [isi] & [isi] \\
		2 layer & [isi] & [isi] & [isi] \\
		3 layer & [isi] & [isi] & [isi] \\
		\bottomrule
	\end{tabular}
	\label{tab:lstm-var-depth}
\end{table}

\textbf{Grafik.} Training loss, validation loss, dan BLEU-4 vs jumlah layer.

\textbf{Analisis.} Bandingkan pengaruh kedalaman pada RNN vs LSTM, termasuk kemampuan menangani long-term dependencies.

\subsection{Perbandingan Ukuran Hidden State: RNN}

\begin{table}[H]
	\centering
	\caption{Variasi Hidden State pada Decoder RNN}
	\begin{tabular}{lccc}
		\toprule
		\textbf{Hidden Units} & \textbf{BLEU-4} & \textbf{METEOR} & \textbf{Waktu Inferensi} \\
		\midrule
		128 & [isi] & [isi] & [isi] \\
		256 & [isi] & [isi] & [isi] \\
		512 & [isi] & [isi] & [isi] \\
		\bottomrule
	\end{tabular}
	\label{tab:rnn-var-hidden}
\end{table}

\textbf{Grafik.} BLEU-4 vs hidden units, serta kurva training/validation loss.

\textbf{Analisis.} Jelaskan pengaruh dimensi hidden state terhadap kapabilitas representasi dan risiko overfitting pada RNN.

\subsection{Perbandingan Ukuran Hidden State: LSTM}

\begin{table}[H]
	\centering
	\caption{Variasi Hidden State pada Decoder LSTM}
	\begin{tabular}{lccc}
		\toprule
		\textbf{Hidden Units} & \textbf{BLEU-4} & \textbf{METEOR} & \textbf{Waktu Inferensi} \\
		\midrule
		128 & [isi] & [isi] & [isi] \\
		256 & [isi] & [isi] & [isi] \\
		512 & [isi] & [isi] & [isi] \\
		\bottomrule
	\end{tabular}
	\label{tab:lstm-var-hidden}
\end{table}

\textbf{Grafik.} BLEU-4 vs hidden units, serta kurva training/validation loss.

\textbf{Analisis.} Bandingkan efek hidden state dimensionality pada RNN vs LSTM, termasuk gradient flow stability.

\subsection{Perbandingan RNN vs LSTM}

\begin{table}[H]
	\centering
	\caption{Perbandingan Decoder Terbaik RNN vs LSTM}
	\begin{tabular}{lccc}
		\toprule
		\textbf{Decoder} & \textbf{BLEU-4} & \textbf{METEOR} & \textbf{Waktu Inferensi} \\
		\midrule
		RNN (best config) & [isi] & [isi] & [isi] \\
		LSTM (best config) & [isi] & [isi] & [isi] \\
		\bottomrule
	\end{tabular}
	\label{tab:rnn-vs-lstm}
\end{table}

\textbf{Grafik.} Kurva training loss, validation loss, dan BLEU-4/METEOR untuk RNN vs LSTM.

\textbf{Analisis kualitatif (minimal 10 contoh).} Sertakan: gambar, caption ground truth, prediksi RNN, prediksi LSTM, serta skor BLEU-4 masing-masing. Kelompokkan contoh menjadi: (1) kasus RNN lebih baik, (2) kasus LSTM lebih baik, (3) kasus keduanya mirip, (4) kasus keduanya buruk.

\subsection{Perbandingan Keras vs From Scratch}

\begin{table}[H]
	\centering
	\caption{Keras vs From Scratch untuk Decoder Terbaik}
	\begin{tabular}{lcccc}
		\toprule
		\textbf{Decoder} & \textbf{Implementasi} & \textbf{BLEU-4} & \textbf{METEOR} & \textbf{Waktu Inferensi} \\
		\midrule
		RNN & Keras & [isi] & [isi] & [isi] \\
		RNN & Scratch & [isi] & [isi] & [isi] \\
		LSTM & Keras & [isi] & [isi] & [isi] \\
		LSTM & Scratch & [isi] & [isi] & [isi] \\
		\bottomrule
	\end{tabular}
	\label{tab:caption-keras-vs-scratch}
\end{table}

\textbf{Grafik.} Kurva training/validation loss untuk Keras dan from-scratch RNN dan LSTM.

\textbf{Analisis.} Verifikasi keselarasan numerik implementasi \textit{from scratch} dengan Keras pada captioning, khususnya pada decoding quality (BLEU-4/METEOR).

\subsection{Pengaruh Panjang Maksimum Caption terhadap BLEU-4}

\begin{table}[H]
	\centering
	\caption{Variasi Panjang Maksimum Caption}
	\begin{tabular}{lcc}
		\toprule
		\textbf{Max Length} & \textbf{BLEU-4} & \textbf{Catatan} \\
		\midrule
		Variasi 1 & [isi] & [isi] \\
		Variasi 2 & [isi] & [isi] \\
		Variasi 3 & [isi] & [isi] \\
		\bottomrule
	\end{tabular}
	\label{tab:caption-maxlen}
\end{table}

\textbf{Grafik.} BLEU-4 dan METEOR vs panjang maksimum caption.

\textbf{Analisis.} Diskusikan trade-off antara panjang caption terhadap kualitas (BLEU-4), kelengkapan informasi, dan computational cost.

\section{Hasil Bonus}

\subsection{Init-Inject vs Pre-Inject}

\begin{table}[H]
	\centering
	\caption{Perbandingan Init-Inject vs Pre-Inject pada Captioning}
	\begin{tabular}{lcccc}
		\toprule
		\textbf{Skema} & \textbf{Merge Mode} & \textbf{BLEU-4} & \textbf{METEOR} & \textbf{Waktu Inferensi} \\
		\midrule
		Pre-Inject & -- & [isi] & [isi] & [isi] \\
		Init-Inject & Add & [isi] & [isi] & [isi] \\
		Init-Inject & Concat & [isi] & [isi] & [isi] \\
		\bottomrule
	\end{tabular}
	\label{tab:init-vs-preinject}
\end{table}

\textbf{Analisis.} Bandingkan kualitas caption dan efisiensi komputasi antara pre-inject (utama) dan init-inject (bonus) dengan dua merge mode, serta dampaknya terhadap gradient flow dan model convergence.

\subsection{Greedy vs Beam Search}

\begin{table}[H]
	\centering
	\caption{Perbandingan Greedy Decoding vs Beam Search}
	\begin{tabular}{lccccc}
		\toprule
		\textbf{Decoder} & \textbf{Strategi} & \textbf{Beam Width} & \textbf{BLEU-4} & \textbf{METEOR} & \textbf{Waktu Inferensi} \\
		\midrule
		RNN & Greedy & -- & [isi] & [isi] & [isi] \\
		RNN & Beam Search & 3 & [isi] & [isi] & [isi] \\
		RNN & Beam Search & 5 & [isi] & [isi] & [isi] \\
		LSTM & Greedy & -- & [isi] & [isi] & [isi] \\
		LSTM & Beam Search & 3 & [isi] & [isi] & [isi] \\
		LSTM & Beam Search & 5 & [isi] & [isi] & [isi] \\
		\bottomrule
	\end{tabular}
	\label{tab:greedy-vs-beam}
\end{table}

\textbf{Analisis.} Analisis trade-off antara kualitas caption (BLEU-4/METEOR) dan waktu komputasi untuk greedy decoding vs beam search dengan berbagai beam width ($k=3, 5$).

\subsection{Visualisasi CNN (Feature Map dan Grad-CAM)}

Sertakan visualisasi untuk arsitektur CNN terbaik:
\begin{enumerate}
	\item \textbf{Feature map intermediate}: Tampilkan aktivasi feature map di beberapa layer (layer 1, tengah, layer terakhir) untuk 3-5 contoh gambar.
	\item \textbf{Grad-CAM overlay}: Overlay heatmap Grad-CAM di atas gambar asli untuk 5 contoh:
	\begin{itemize}
		\item 2 contoh klasifikasi benar (high confidence).
		\item 2 contoh klasifikasi salah (model prediksi lain).
		\item 1 contoh borderline (low confidence).
	\end{itemize}
	\item \textbf{Analisis visual}: Diskusikan apakah heatmap selaras dengan region penting dalam gambar (e.g., objek utama).
\end{enumerate}

\subsection{Backward Propagation From Scratch}

\begin{table}[H]
	\centering
	\caption{Verifikasi Backward Propagation pada Layer Kritikal}
	\begin{tabular}{lcc}
		\toprule
		\textbf{Layer} & \textbf{Gradient Check (Numerik vs Analitik)} & \textbf{Status} \\
		\midrule
		Conv2D & [isi] & [Passed/Failed] \\
		MaxPool2D & [isi] & [Passed/Failed] \\
		AvgPool2D & [isi] & [Passed/Failed] \\
		Dense & [isi] & [Passed/Failed] \\
		LSTMCell & [isi] & [Passed/Failed] \\
		\bottomrule
	\end{tabular}
	\label{tab:gradient-check}
\end{table}

\textbf{Verifikasi numeric.} Sertakan:
\begin{itemize}
	\item Gradient checking dengan finite difference (h = 1e-5) untuk setiap layer.
	\item Visualisasi: plot perbedaan relative antara numerik dan analitik gradient.
	\item Ringkasan: apakah semua gradient check pass (relative error < 1e-4)?
\end{itemize}

\section{Ringkasan Temuan}

Tuliskan poin-poin temuan utama setelah seluruh tabel dan grafik diisi, dengan struktur berikut:

\subsection{CNN Classification}

\begin{enumerate}
	\item \textbf{Konfigurasi terbaik}: Laporan arsitektur CNN dengan macro F1-score tertinggi, termasuk jumlah layer, filter distribution, kernel size, dan pooling type.
	\item \textbf{Pengaruh hyperparameter}:
	\begin{itemize}
		\item Depth: apakah lebih dalam selalu lebih baik? Adakah tanda overfitting?
		\item Filter count: trade-off antara model capacity dan accuracy.
		\item Kernel size: pengaruh receptive field terhadap spatial pattern recognition.
		\item Pooling strategy: Max vs Average dalam konteks dataset Intel Image.
	\end{itemize}
	\item \textbf{Shared vs non-shared}: Analisis keuntungan/kerugian Conv2D vs LocallyConnected2D dalam hal akurasi, parameter efficiency, dan training time.
	\item \textbf{From scratch vs Keras}: Verifikasi keselarasan numerik dan gradient correctness.
\end{enumerate}

\subsection{Image Captioning (RNN/LSTM)}

\begin{enumerate}
	\item \textbf{Konfigurasi terbaik}: Laporan decoder terbaik (RNN atau LSTM) dengan BLEU-4 tertinggi, termasuk jumlah layer dan hidden units.
	\item \textbf{RNN vs LSTM}: Analisis perbedaan dalam kemampuan menangani long-term dependencies dan kualitas caption.
	\item \textbf{Pengaruh hyperparameter pada captioning}:
	\begin{itemize}
		\item Depth: pengaruh pada representasi kontekstual dan gradient flow.
		\item Hidden units: trade-off antara kapabilitas dan overfitting.
		\item Max caption length: dampak terhadap BLEU-4 dan computational cost.
	\end{itemize}
	\item \textbf{Decoding strategy}: Analisis greedy vs beam search ($k=3, 5$) dalam kualitas dan efisiensi.
	\item \textbf{Bonus features}: Perbandingan init-inject vs pre-inject dan pengaruh merge mode ("add" vs "concat").
	\item \textbf{From scratch vs Keras}: Verifikasi keselarasan pada BLEU-4 dan loss curve.
\end{enumerate}

\subsection{Kesimpulan Umum}

\begin{enumerate}
	\item Insight utama tentang design choices yang optimal untuk kedua tugas.
	\item Perbedaan karakteristik CNN (discriminative) vs captioning (generative).
	\item Rekomendasi best practices untuk implementasi \textit{from scratch} yang sesuai Keras.
\end{enumerate}
